\documentclass[11pt]{article}
\usepackage[margin=1in]{geometry}
\usepackage{amsmath,amssymb,amsthm,mathtools}
\usepackage{bbm}
\usepackage{enumitem}
\usepackage{hyperref}

\newtheorem{definition}{Definition}[section]
\newtheorem{theorem}{Theorem}[section]
\newtheorem{proposition}{Proposition}[section]
\newtheorem{lemma}{Lemma}[section]
\newtheorem{corollary}{Corollary}[section]
\theoremstyle{remark}
\newtheorem{remark}{Remark}[section]
\newtheorem{example}{Example}[section]

\newcommand{\R}{\mathbb{R}}
\newcommand{\N}{\mathbb{N}}
\newcommand{\PP}{\mathbb{P}}
\newcommand{\EE}{\mathbb{E}}
\newcommand{\1}{\mathbbm{1}}
\newcommand{\cF}{\mathcal{F}}
\newcommand{\cB}{\mathcal{B}}
\newcommand{\cA}{\mathcal{A}}
\newcommand{\cC}{\mathcal{C}}

\title{Probability Foundations for Econometrics\\\large An Intuition-First Rewrite of the NYU Metrics Notes}
\author{Adapted for advanced master's students in economics}
\date{\today}

\begin{document}

\maketitle

These notes rewrite the measure-theoretic probability material behind econometrics in more direct language. The goal is not to remove rigor, but to make clear what each object means, why we need it, and where it enters econometric arguments.

\tableofcontents

\section{How to Read These Notes}

If you come from economics rather than pure mathematics, the hardest part of probability theory is often the notation, not the ideas. Three principles help:

\begin{enumerate}[label=(\arabic*)]
  \item A probability model needs a set of outcomes, a collection of events we are allowed to talk about, and a rule that assigns probabilities to those events.
  \item A random variable is just a measurable way of translating outcomes into numbers.
  \item Expectations, laws of large numbers, and asymptotic results are all built on integration.
\end{enumerate}

You do not need to love set theory to use these notes. What matters is understanding what the objects let us do in econometrics: define distributions, compute expectations, state convergence results, and justify asymptotic approximations.

\section{Probability Spaces}

\subsection{The three ingredients}

\begin{definition}[Probability space]
A \emph{probability space} is a triple \((\Omega,\cF,\PP)\), where:
\begin{enumerate}[label=(\roman*)]
  \item \(\Omega\) is the set of all possible outcomes;
  \item \(\cF\) is a collection of subsets of \(\Omega\), called \emph{events};
  \item \(\PP\) assigns a probability to each event in \(\cF\).
\end{enumerate}
\end{definition}

\begin{remark}
Think of \(\Omega\) as the full hidden world, \(\cF\) as the events we are allowed to ask probability questions about, and \(\PP\) as the model's answer.
\end{remark}

\begin{example}[Simple economic example]
Suppose one person is drawn from a population. An outcome \(\omega\) could encode that person's income, education, gender, region, and treatment status. Then an event might be
\[
  A = \{\omega : \text{income}(\omega) > 50{,}000\}.
\]
Probability theory lets us assign a number \(\PP(A)\) to such an event.
\end{example}

\subsection{Why we need a sigma-algebra}

\begin{definition}[\(\sigma\)-algebra]
A collection \(\cF\) of subsets of \(\Omega\) is a \emph{\(\sigma\)-algebra} if:
\begin{enumerate}[label=(\roman*)]
  \item \(\emptyset \in \cF\) and \(\Omega \in \cF\);
  \item if \(A \in \cF\), then \(A^c \in \cF\);
  \item if \(A_1,A_2,\dots \in \cF\), then \(\bigcup_{i=1}^{\infty} A_i \in \cF\).
\end{enumerate}
\end{definition}

\begin{remark}
This definition looks technical, but the idea is simple: once an event is meaningful, its complement should also be meaningful, and countable combinations of meaningful events should remain meaningful.
\end{remark}

\begin{example}[Treatment assignment]
Suppose \(\Omega\) records whether a unit is treated and whether it has high income. If we can talk about the event ``treated,'' we should also be able to talk about ``not treated,'' ``treated and high income,'' and ``treated or high income.'' A \(\sigma\)-algebra guarantees this closure.
\end{example}

\begin{definition}[Algebra]
A collection \(\cA\) of subsets of \(\Omega\) is an \emph{algebra} if it is closed under complements and under \emph{finite} unions.
\end{definition}

\begin{remark}
So an algebra handles only finite combinations. A \(\sigma\)-algebra is stronger because it also handles countably many operations. That extra strength matters for limits, asymptotics, and convergence arguments.
\end{remark}

\begin{definition}[Probability measure]
Given a measurable space \((\Omega,\cF)\), a function \(\PP : \cF \to [0,1]\) is a \emph{probability measure} if:
\begin{enumerate}[label=(\roman*)]
  \item \(\PP(A) \ge 0\) for every \(A \in \cF\);
  \item \(\PP(\Omega) = 1\);
  \item if \(A_1,A_2,\dots\) are pairwise disjoint events, then
  \[
    \PP\Bigl(\bigcup_{i=1}^{\infty} A_i\Bigr) = \sum_{i=1}^{\infty}\PP(A_i).
  \]
\end{enumerate}
\end{definition}

\begin{remark}
The third property is countable additivity. It is the key rule that makes probability compatible with infinite sequences and limiting arguments.
\end{remark}

\subsection{Basic properties you use all the time}

Once \((\Omega,\cF,\PP)\) is given, the following facts follow immediately:

\begin{enumerate}[label=(\alph*)]
  \item \(\PP(\emptyset)=0\).
  \item \(\PP(A^c)=1-\PP(A)\).
  \item If \(A \subseteq B\), then \(\PP(A)\le \PP(B)\).
  \item \(\PP(A\cup B)=\PP(A)+\PP(B)-\PP(A\cap B)\).
  \item If \(A_n \uparrow A\), then \(\PP(A_n)\uparrow \PP(A)\).
  \item If \(A_n \downarrow A\), then \(\PP(A_n)\downarrow \PP(A)\).
  \item For any events \(A_1,A_2,\dots\),
  \[
    \PP\Bigl(\bigcup_{i=1}^{\infty} A_i\Bigr)\le \sum_{i=1}^{\infty}\PP(A_i).
  \]
\end{enumerate}

\begin{remark}[Why economists care]
Statements like ``the probability of a large estimation error goes to zero'' are really statements about probabilities of sequences of events. The continuity properties above are what make those statements mathematically sound.
\end{remark}

\section{Generated Sigma-Algebras and Borel Sets}

\subsection{Generated sigma-algebras}

\begin{definition}[Generated \(\sigma\)-algebra]
If \(\cC\) is any collection of subsets of \(\Omega\), then the \emph{\(\sigma\)-algebra generated by \(\cC\)} is
\[
  \sigma(\cC) := \bigcap \left\{\cF : \cF \text{ is a \(\sigma\)-algebra and } \cC \subseteq \cF\right\}.
\]
\end{definition}

\begin{remark}
\(\sigma(\cC)\) is the smallest \(\sigma\)-algebra containing the sets in \(\cC\). Intuitively, you start from a few primitive events and then add everything needed to make the collection closed under complements and countable unions.
\end{remark}

\begin{example}[A dice example]
Let \(\Omega=\{1,2,3,4,5,6\}\) and suppose we only care whether the outcome is in \(\{1,2\}\). If \(\cC=\{\{1,2\}\}\), then
\[
  \sigma(\cC)=\{\emptyset,\{1,2\},\{3,4,5,6\},\Omega\}.
\]
So this information structure only distinguishes ``1 or 2'' from ``anything else.''
\end{example}

\subsection{The Borel sigma-algebra}

\begin{definition}[Borel \(\sigma\)-algebra on \(\R\)]
The \emph{Borel \(\sigma\)-algebra} on \(\R\), denoted \(\cB\), is the \(\sigma\)-algebra generated by open intervals:
\[
  \cB = \sigma\bigl(\{(a,b): a<b\}\bigr).
\]
\end{definition}

\begin{remark}
In practice, Borel sets are the usual sets you want to assign probabilities to: intervals, unions of intervals, rectangles in \(\R^k\), and the sets built from them by countably many operations.
\end{remark}

On \(\R^k\), we write \(\cB^k\) for the Borel \(\sigma\)-algebra generated by rectangles such as
\[
  (-\infty,a_1]\times \cdots \times (-\infty,a_k].
\]

\begin{remark}[Why this matters]
If \(X\) is income, \(Y\) is education, and \(D\) is treatment, then events like
\[
  \{X \le x\}, \qquad \{X \le x,\ Y \le y\}
\]
need to be measurable. The Borel \(\sigma\)-algebra is the standard way to formalize that.
\end{remark}

\subsection{Lebesgue measure and outer measure}

On \(\R\), the most natural notion of size is interval length. Lebesgue measure extends that idea from intervals to a very large class of sets.

\begin{definition}[Lebesgue measure, informal]
Lebesgue measure assigns to an interval its length and extends this consistently to a large \(\sigma\)-algebra of subsets of \(\R\).
\end{definition}

\begin{remark}
For econometrics, Lebesgue measure is important because continuous distributions are defined relative to it. Saying that a random variable has a density means that its distribution can be represented using ordinary integration with respect to length, area, volume, and so on.
\end{remark}

\begin{definition}[Outer measure]
Given a non-negative set function \(\mu\) defined on a simple class of sets, the \emph{outer measure} of \(A\) is
\[
  \mu^*(A) = \inf\left\{\sum_{j=1}^{\infty}\mu(A_j) : A \subseteq \bigcup_{j=1}^{\infty} A_j\right\},
\]
where the \(A_j\)'s come from the simple class.
\end{definition}

\begin{remark}
You will almost never compute outer measures in applied econometrics. The point is conceptual: this construction explains how one goes from simple objects such as intervals to a rigorous measure defined on a rich enough collection of sets.
\end{remark}

\section{Random Variables and Their Distributions}

\subsection{Random variables as measurable functions}

\begin{definition}[Random variable and random vector]
Let \((\Omega,\cF,\PP)\) be a probability space. A function \(X:\Omega \to \R^k\) is a \emph{random vector} if for every \(B \in \cB^k\),
\[
  X^{-1}(B)=\{\omega \in \Omega : X(\omega)\in B\}\in \cF.
\]
When \(k=1\), we call \(X\) a \emph{random variable}.
\end{definition}

\begin{remark}
This is the definition of measurability. It ensures that statements such as ``\(X \le x\)'' or ``\(X\) belongs to a set \(B\)'' are genuine events to which probabilities can be assigned.
\end{remark}

\begin{example}[Regression language]
If \(Y_i\) is earnings and \(X_i\) is schooling, then events like \(\{Y_i \le 30{,}000\}\) or \(\{Y_i \le 30{,}000,\ X_i \ge 16\}\) must be measurable. Otherwise the expressions \(\PP(Y_i \le 30{,}000)\) and \(\PP(Y_i \le 30{,}000,\ X_i \ge 16)\) would not even be well-defined.
\end{example}

\subsection{Distribution and induced probability measure}

\begin{definition}[Distribution of a random vector]
If \(X:\Omega \to \R^k\) is a random vector, its \emph{distribution} or \emph{law} is the probability measure \(P_X\) on \((\R^k,\cB^k)\) defined by
\[
  P_X(B)=\PP(X\in B), \qquad B \in \cB^k.
\]
\end{definition}

\begin{remark}
This definition says that the random variable \(X\) transfers the original probability measure \(\PP\) on the abstract space \(\Omega\) onto a probability measure on \(\R^k\). Once we know \(P_X\), we can usually work directly on \(\R^k\) and stop worrying about the hidden underlying sample space.
\end{remark}

\begin{definition}[Cumulative distribution function]
If \(X\) is real-valued, its cumulative distribution function (CDF) is
\[
  F_X(x)=\PP(X\le x).
\]
If \(X=(X_1,\dots,X_k)\in \R^k\), then
\[
  F_X(x_1,\dots,x_k)=\PP(X_1\le x_1,\dots,X_k\le x_k).
\]
\end{definition}

\begin{proposition}[Basic properties of a CDF]
For a real-valued random variable \(X\), the function \(F_X\) is:
\begin{enumerate}[label=(\roman*)]
  \item non-decreasing;
  \item right-continuous;
  \item such that \(\lim_{x\to -\infty}F_X(x)=0\) and \(\lim_{x\to \infty}F_X(x)=1\).
\end{enumerate}
\end{proposition}

\begin{remark}
The CDF contains the same information as the distribution. In one dimension, knowing \(F_X\) is enough to recover probabilities of intervals and other Borel sets.
\end{remark}

\subsection{The information contained in a random variable}

\begin{definition}[\(\sigma\)-algebra generated by \(X\)]
The \(\sigma\)-algebra generated by a random variable \(X\) is
\[
  \sigma(X)=\{X^{-1}(B): B \in \cB\}.
\]
\end{definition}

\begin{remark}
\(\sigma(X)\) is the collection of all events that can be described using \(X\) alone. So it is natural to interpret \(\sigma(X)\) as the information revealed by \(X\).
\end{remark}

\begin{example}[Information comparison]
If \(Y=2X+1\), then \(Y\) contains exactly the same information as \(X\), because knowing one tells you the other. In that case \(\sigma(Y)=\sigma(X)\). By contrast, if \(Y=\1\{X>0\}\), then \(Y\) reveals less information, so \(\sigma(Y)\subseteq \sigma(X)\).
\end{example}

\section{Absolute Continuity and Densities}

\subsection{What a density really means}

Let \(\lambda^k\) denote Lebesgue measure on \(\R^k\).

\begin{definition}[Absolute continuity]
A probability measure \(P_X\) on \((\R^k,\cB^k)\) is \emph{absolutely continuous} with respect to Lebesgue measure, written \(P_X \ll \lambda^k\), if
\[
  \lambda^k(B)=0 \quad \Rightarrow \quad P_X(B)=0
\]
for every Borel set \(B\).
\end{definition}

\begin{remark}
The meaning is: any set that has zero geometric size also has zero probability. This is what we expect from a genuinely continuous distribution.
\end{remark}

\begin{theorem}[Existence of a density]
If \(P_X \ll \lambda^k\), then there exists a non-negative measurable function \(f_X\) such that
\[
  P_X(B)=\int_B f_X(u)\,d\lambda^k(u)
\]
for every Borel set \(B\). The function \(f_X\) is called a \emph{density} of \(X\).
\end{theorem}

\begin{remark}
This theorem is the formal reason why we are allowed to write familiar expressions like
\[
  \PP(a \le X \le b)=\int_a^b f_X(x)\,dx.
\]
\end{remark}

If \(X\) is one-dimensional and has density \(f_X\), then
\[
  F_X(x)=\int_{-\infty}^x f_X(t)\,dt.
\]

\begin{example}[Continuous versus discrete]
If \(X \sim N(0,1)\), then \(X\) has a density. If \(X\) is Bernoulli, it does not have a density with respect to Lebesgue measure on \(\R\); instead, it puts point mass on \(0\) and \(1\).
\end{example}

\section{Integration and Expectation}

\subsection{Why integration enters probability}

In econometrics, we often care about quantities such as
\[
  \EE[X], \qquad \EE[g(X)], \qquad \EE[X \mid Z], \qquad \EE[m(W,\theta)].
\]
All of these are integrals. So to understand probability rigorously, we need a theory of integration that works for general random variables.

\subsection{Measurable functions and simple functions}

\begin{definition}[Measurable function]
A function \(g:\R^k \to \R\) is measurable if \(g^{-1}(B)\in \cB^k\) for every Borel set \(B\in \cB\).
\end{definition}

\begin{remark}
In practice, the functions used in econometrics are almost always measurable: polynomials, indicators, absolute values, exponentials, likelihood contributions, moments, and so on.
\end{remark}

\begin{definition}[Simple function]
A \emph{simple function} is a finite step function of the form
\[
  s(u)=\sum_{i=1}^{n} a_i \1_{B_i}(u),
\]
where the \(a_i\)'s are real numbers and the \(B_i\)'s are measurable sets.
\end{definition}

\begin{remark}
Simple functions are important because the Lebesgue integral is built by approximating more complicated functions with step functions.
\end{remark}

\subsection{Expectation as an integral}

\begin{definition}[Integral of a simple function]
If \(s(u)=\sum_{i=1}^{n} a_i \1_{B_i}(u)\), then
\[
  \int s(u)\,dP_X(u)=\sum_{i=1}^{n} a_i P_X(B_i).
\]
\end{definition}

From simple functions one defines the integral of non-negative measurable functions, and then of general integrable functions.

\begin{definition}[Expectation]
If \(X\) has distribution \(P_X\) and \(g\) is integrable, then
\[
  \EE[g(X)] = \int_{\R^k} g(u)\,dP_X(u).
\]
\end{definition}

\begin{remark}[Three useful special cases]
\begin{enumerate}[label=(\roman*)]
  \item If \(X\) is discrete with support \(\{x_j\}\), then
  \[
    \EE[g(X)] = \sum_j g(x_j)\PP(X=x_j).
  \]
  \item If \(X\) has density \(f_X\), then
  \[
    \EE[g(X)] = \int_{\R^k} g(u)f_X(u)\,du.
  \]
  \item If \(A\) is an event, then
  \[
    \EE[\1_A]=\PP(A).
  \]
\end{enumerate}
\end{remark}

\begin{proposition}[Basic properties of expectation]
Whenever the expectations are well-defined:
\begin{enumerate}[label=(\alph*)]
  \item If \(g_1 \le g_2\), then \(\EE[g_1(X)] \le \EE[g_2(X)]\).
  \item \(\EE[a g_1(X) + b g_2(X)] = a\EE[g_1(X)] + b\EE[g_2(X)]\).
  \item \(|\EE[g(X)]| \le \EE[|g(X)|]\).
\end{enumerate}
\end{proposition}

\begin{remark}[LOTUS]
The formula
\[
  \EE[g(X)] = \int g(u)\,dP_X(u)
\]
is sometimes called the \emph{law of the unconscious statistician}. It says that once we know the distribution of \(X\), we can compute expectations without going back to the original sample space.
\end{remark}

\subsection{Integration on the abstract probability space}

We can also define expectation directly on \((\Omega,\cF,\PP)\). For a simple random variable
\[
  S(\omega)=\sum_{i=1}^{m} a_i \1_{B_i}(\omega),
\]
we set
\[
  \int S(\omega)\,d\PP(\omega)=\sum_{i=1}^{m} a_i \PP(B_i).
\]

\begin{remark}
Whether we integrate over the abstract space \((\Omega,\cF,\PP)\) or over the induced distribution \((\R^k,\cB^k,P_X)\), we are describing the same expectation. The second viewpoint is often easier in applications.
\end{remark}

\section{Convergence and the Main Limit Theorems}

\subsection{Why convergence matters}

Asymptotic econometrics studies what happens when sample size \(n\) becomes large. To do that, we need precise notions of convergence. Different notions are useful for different results.

\subsection{Modes of convergence}

Let \((X_n)_{n\ge 1}\) and \(X\) be random variables on the same probability space.

\begin{definition}[Almost sure convergence]
We say \(X_n \to X\) \emph{almost surely} if
\[
  \PP\left(\left\{\omega : X_n(\omega) \to X(\omega)\right\}\right)=1.
\]
\end{definition}

\begin{definition}[Convergence in probability]
We say \(X_n \xrightarrow{p} X\) if for every \(\varepsilon>0\),
\[
  \PP(|X_n-X|>\varepsilon)\to 0.
\]
\end{definition}

\begin{definition}[Convergence in \(L^1\)]
We say \(X_n \to X\) in \(L^1\) if
\[
  \EE[|X_n-X|]\to 0.
\]
\end{definition}

\begin{definition}[Convergence in distribution]
We say \(X_n \xrightarrow{d} X\) if
\[
  F_{X_n}(x)\to F_X(x)
\]
for every continuity point \(x\) of \(F_X\).
\end{definition}

\begin{remark}[How to think about them]
\begin{enumerate}[label=(\roman*)]
  \item Almost sure convergence is path-by-path convergence, except on a probability-zero set.
  \item Convergence in probability means large deviations become unlikely.
  \item \(L^1\) convergence means the expected absolute error vanishes.
  \item Convergence in distribution means the distributions get closer, even if the variables are not defined on the same sample paths in a useful way.
\end{enumerate}
\end{remark}

\begin{proposition}[Standard implications]
We have
\[
  X_n \to X \text{ a.s. } \Rightarrow X_n \xrightarrow{p} X,
\]
and
\[
  X_n \xrightarrow{p} X \Rightarrow X_n \xrightarrow{d} X.
\]
Also,
\[
  X_n \to X \text{ in } L^1 \Rightarrow X_n \xrightarrow{p} X.
\]
\end{proposition}

\subsection{Monotone convergence theorem}

\begin{theorem}[MCT]
Let \(g_n:\R^k \to [0,\infty]\) be measurable and suppose
\[
  g_n(u)\uparrow g(u)
\]
for every \(u\). Then
\[
  \int g_n(u)\,dP_X(u) \to \int g(u)\,dP_X(u).
\]
\end{theorem}

\begin{remark}
If the functions increase steadily to a limit and never become negative, then we may pass the limit inside the integral. No dominating bound is needed.
\end{remark}

\subsection{Dominated convergence theorem}

\begin{theorem}[DCT]
Suppose \(g_n:\R^k \to \R\) are measurable, \(g_n(u)\to g(u)\) pointwise, and there exists an integrable function \(M\) such that
\[
  |g_n(u)| \le M(u)
\]
for every \(u\) and every \(n\). Then
\[
  \int g_n(u)\,dP_X(u) \to \int g(u)\,dP_X(u).
\]
\end{theorem}

\begin{remark}[Why DCT matters in econometrics]
DCT is the theorem behind many familiar operations: exchanging limits and expectations, proving consistency of sample moments, and showing that criterion functions converge. The crucial requirement is a single integrable bound \(M\) that dominates the entire sequence.
\end{remark}

\section{Conditional Probability and Independence}

\subsection{Conditional probability}

\begin{definition}[Conditional probability]
If \(A,B \in \cF\) and \(\PP(B)>0\), then
\[
  \PP(A \mid B) = \frac{\PP(A \cap B)}{\PP(B)}.
\]
\end{definition}

\begin{remark}
Conditioning means we restrict attention to the subpopulation described by \(B\). The denominator rescales probabilities so that the conditional world has total probability one.
\end{remark}

\begin{proposition}[Bayes' rule]
If \(\PP(A)>0\) and \(\PP(B)>0\), then
\[
  \PP(A \mid B)=\frac{\PP(B \mid A)\PP(A)}{\PP(B)}.
\]
\end{proposition}

\begin{remark}
Bayes' rule is often useful when \(\PP(B \mid A)\) is easier to model than \(\PP(A \mid B)\).
\end{remark}

\subsection{Independence}

\begin{definition}[Independence of events]
Events \(A_1,\dots,A_k\) are independent if for every subcollection,
\[
  \PP(A_{i_1}\cap \cdots \cap A_{i_m}) = \PP(A_{i_1})\cdots \PP(A_{i_m}).
\]
\end{definition}

\begin{definition}[Independence of random variables]
Random variables \(X_1,\dots,X_k\) are independent if for all Borel sets \(B_1,\dots,B_k\),
\[
  \PP(X_1 \in B_1,\dots,X_k \in B_k)=\prod_{j=1}^{k}\PP(X_j \in B_j).
\]
\end{definition}

\begin{remark}
Independence is much stronger than zero correlation. Independence means the joint distribution factorizes completely; zero correlation is only one moment condition.
\end{remark}

For two real-valued random variables, independence implies
\[
  F_{(X_1,X_2)}(x_1,x_2)=F_{X_1}(x_1)F_{X_2}(x_2).
\]

\begin{remark}[Econometric interpretation]
When we assume i.i.d.\ sampling, we are making two claims at once:
\begin{enumerate}[label=(\roman*)]
  \item independence across observations;
  \item a common marginal distribution.
\end{enumerate}
Both parts matter for large-sample theory.
\end{remark}

\section{Law of Large Numbers}

\subsection{Statement}

\begin{theorem}[Weak law of large numbers]
Let \(X_1,X_2,\dots\) be i.i.d.\ random variables with \(\EE[|X_1|]<\infty\). Define
\[
  \mu=\EE[X_1], \qquad \bar X_n = \frac{1}{n}\sum_{i=1}^{n} X_i.
\]
Then
\[
  \bar X_n \xrightarrow{p} \mu.
\]
\end{theorem}

\begin{remark}
This is the foundational consistency result behind sample averages, sample moments, many estimators defined by moment conditions, and the behavior of empirical objective functions.
\end{remark}

\begin{example}[Sample mean]
If \(Y_1,\dots,Y_n\) are i.i.d.\ wages with finite mean, then the sample average wage converges in probability to the population mean wage.
\end{example}

\subsection{Why the assumptions matter}

\begin{enumerate}[label=(\alph*)]
  \item \emph{Independence} prevents observations from carrying the same information repeatedly.
  \item \emph{Identical distribution} ensures there is one stable target parameter \(\mu\).
  \item \emph{Finite first moment} guarantees the mean is well-defined.
\end{enumerate}

\begin{remark}
In more advanced econometrics, these assumptions are weakened in many directions: dependence over time, heterogeneity across units, triangular arrays, mixing processes, and martingale difference structures. But the logic remains the same: convergence of averages requires probabilistic control of dependence and tail behavior.
\end{remark}

\section{Final Perspective}

The main takeaway is that measure-theoretic probability is not a separate subject from econometrics. It is the language that lets us state and prove the results economists use all the time:

\begin{enumerate}[label=(\roman*)]
  \item random variables have distributions;
  \item expectations are integrals;
  \item convergence has several distinct meanings;
  \item asymptotic theorems such as the LLN rely on those definitions.
\end{enumerate}

Once these foundations are clear, later topics such as the central limit theorem, conditional expectation, likelihood methods, GMM, and modern asymptotic theory become much easier to organize conceptually.

\end{document}
